\documentclass[12pt,a4paper]{article}
\usepackage[scheme=plain]{ctex}
\usepackage[margin=2.5cm]{geometry}
\usepackage{indentfirst}
\usepackage{setspace}
\usepackage{fontspec}
\setCJKmainfont[BoldFont=SimHei,ItalicFont=KaiTi]{SimSun}
\newCJKfontfamily\kaifont{KaiTi}[AutoFakeBold]
\usepackage{newpxtext,newpxmath}

\usepackage{lastpage}





\pagestyle{plain}

\usepackage[nobottomtitles*]{titlesec}
\titleformat{\section}{\Large\bfseries}{}{0em}{ }
\titlespacing*{\section}{0pt}{3ex}{1.5ex}

\usepackage{xcolor}
\usepackage{needspace}

\setlength{\parindent}{2em}
\setlength{\parskip}{0.4em}

\doublespacing

\title{\bfseries Day 6\\[0.3em]组织的纹理}
\author{}
\date{}

\begin{document}
\maketitle

\vspace{-2em}

\vspace{1em}

\section{前哨的第六封信}
\vspace{-0.3em}{\color{gray}\small\kaifont \noindent K 在村子里待了足够久。白天和村长谈话，晚上在小屋里整理自己的所见。他不再试图从外部观察城堡，也不再追问"它到底在哪"。他把这几天看到的几件事拼在一起，写给即将出发的队友。}\vspace{0.5em}
\vspace{0.3em}

\noindent\rule{\textwidth}{0.4pt}

\vspace{0.3em}



同志们：



你们听完了周雪光，也讨论过 Lipsky——现在我想带你们回到城堡本身。



我在村里看到了三样东西。公文在每个环节都有答复，但从没有裁断：下级负责遗忘，上级负责铭记。职能之间是彻底区隔的，一个部门不可以问另一个部门“是不是你发的”——这是纹理，是组织在生长和断裂中不变的东西。然后巴纳巴斯去了办公厅，穿着自己的衣服，站在那些他能越过的和不能越过的栅栏之间，发现它们看上去一模一样。



这三件事是同一件事。你们先读，读完请告诉我你们看到了什么。



K

\newpage
\section{电话与公文}
\vspace{-0.3em}{\color{gray}\small\kaifont \noindent K 来到村公所，村长卧病在床。他讲了两件事。打到城堡的电话会让最下一级部门的全部电话一起响起来，即便有人接听了也不能当真——那是"八音盒"。公文在各部门之间来回流转，每个环节都有回应，但最终没有裁断。如果说空间上的隔绝和人际关系里的渗透是你能感受到的痕迹，那制度上的错位呢？}\vspace{0.5em}
\vspace{0.3em}

\noindent\rule{\textwidth}{0.4pt}

\vspace{0.3em}



K对此地衙门的这种看法，首先在村长处得到了完全的证实。这是个和颜悦色、胡须刮得干干净净的胖子，他是在病中，在风湿病严重发作时躺在床上接见K的。“这么说您就是我们的土地测量员先生了。”他说，同时想坐起来表示欢迎，但怎么使劲也无用，于是带着歉意指指自己的腿又颓然躺回枕头上去。一个不苟言笑的女人给K拿来了一把椅子摆在床前，这间屋子的窗子本来就很小，又拉上了窗帘而更加昏暗，所以几乎只能看得见这女人那影影绰绰的轮廓。“您请坐，您请坐，土地测量员先生，”村长说，“请告诉我您有什么要求。”K念了克拉姆的信又加上几句自己的话。现在他再次体味到同衙门打交道时那种异常轻松的感觉。衙门简直就肩负着所有的重担，您可以把什么都推给它，自己落得个轻松自在。村长似乎从他那一方面同样感到了这一点，在床上很不舒服地翻了个身。最后他说：“您已经看到了，土地测量员先生，整个这件事情我是早就知道了的，我之所以还没有采取相应措施，第一是因为我生病，另外就是您很久没来，我已经在想您大概已经放弃这工作了。但现在呢，既然您亲临寒舍造访，我当然就必须把令人不快的全部实情倾囊相告。如您所说，您已被录用为土地测量员；但是遗憾得很，我们并不需要土地测量员。这里可以说一点让土地测量员做的工作也没有。我们这些小家小户的土地界线是划定了的，一应事项已全部有条不紊登记在案。产业易主的事极少发生，小的边界纠纷我们自己来解决。所以我们要土地测量员做什么？”虽然K原先完全没有细想过这一点，但他内心深处确信，自己是早就料到会得到类似说法了。〔7〕



正因为如此他立即说道：“我真万万没有想到事情是这样。这使我的全部打算都落空了。现在我只希望这是一个误会。”——“可惜没有误会，”村长说，“情况正像我说的那样。”——“可这也太不像话了！”K叫起来。“难道我千辛万苦长途跋涉到这里就为了让人再把我送回去吗！”——“这是另一码事，”村长说，“这件事我无权决定，但您所谓的那个误会是怎么产生的我倒可以给您解释解释。对一个像伯爵衙门这样的大衙门来说，有时难免会发生一个部门发出这一项指令、另一个部门发出另一项指令而互相不通气的事，虽然上一级的监督检查是极为严格的，但检查很自然地往往晚到一步，于是总会出点小差错，当然毫无例外都是在一些微乎其微的小事上出错，比如您这件事。在大事上我还没听说出过什么错，不过就是小事也够让人尴尬的了。至于说到您这件事，那么我可以坦率地把事情原原本本讲给您听，我用不着保守公务秘密——要保密我的官还太小，我是务农的，一辈子务农。很久以前，那时我当村长才几个月，上头来了一份公函，我记不清是哪个部门发出的了，那公函用高级长官特有的命令口气通知说，拟聘任一名土地测量员，指令村公所将有关该测量员工作必需的各项计划和全部文字材料立即备齐待命。这公函上指的人当然不可能是您，因为这是好多年前的事了，要不是现在我卧病在床，有的是时间去回想那些最最可笑的事情，那么我也是不会记起这事来的。”——“米齐，”他突然中断叙述对那个一直在屋里跑来跑去不知在忙乎什么的女人说，“请你去那边柜里看看，兴许能找到那份文件。”——“这项指令，”他接着对K解释，“是我刚当上村长那段时间下达的，那时我把所有的文件都保存起来了。”那女人当即打开柜子，K和村长在一旁看着她。柜子里各种文件纸张塞得满满当当的。柜门一开，两大捆像劈柴一样被捆成一大卷的文件便骨碌碌滚了出来，把那女人吓了一跳，一下子闪开了。“下面，可能在下面，”村长在床上指挥着，那女人则驯顺地张开双臂把大批大批的文件抱出来撂在地上，这样才好去拿最底下的那份文件。现在一卷卷文牍纸片已经堆满了半间屋子。“我们做了很多工作呵，”村长颔首说，“这些还只是一小部分呢。大宗的我都存放在库房里，而且，绝大部分文件都丢失了。谁有本事把全部文件都保存下来呀！库房里还有很多很多呢。”——“你找得着那份指令吗？”过一会儿他又对他女人说。“你先得找一份上面有‘土地测量员’字样，底下还画了蓝道道的文件。”——“这儿太暗了，”那女人说，“我去取一支蜡烛来，”说罢就迈步跨过地上的大堆文牍出屋去了。“我女人在这项繁重的公务中帮了我的大忙了，”村长说，“我还有别的重要事务，这项工作只能算是兼管一下。虽说我还有一个助手帮助我整理这些文字材料，这就是那位老师，可即使这样也仍然做不完，总是剩下许多堆积起来，还没整理的全部集中在那边那个柜子里，”他指指另一个柜子。“现在我这一病倒呢，那里简直就堆成山了，”说完这句他又疲倦地、然而却面带骄傲之色躺了下去。“我可不可以，”当那女人取来了蜡烛后又跪在柜子前继续找那份指令时，K说道，“帮着您太太一起找？”村长微笑着摇摇头：“刚才我说过，对您我没有什么公务秘密可保；可是，要让您本人去文件堆里找东西，那我确实又走得太远了。”此刻屋里一片安静，唯闻翻动纸张窸窸窣窣之声，在这种情况下村长甚至微微打起盹来。一阵轻轻的叩门声使K回身一看，当然又是两个助手。不过无论如何，两人现在总算有点长进：他们不是乒乒乓乓闯进屋来，而是站在外面通过微开的门缝轻声说：“我们在外面待着太冷了。”——“谁呀？”村长被惊醒了，问道。“没别人，是我的两个助手，”K说，“我不知道该让他们在哪儿等我，外面太冷，这儿嘛他们又让人讨厌。”——“我不在乎他们，”村长和气地说。“您让他们进来好了。其实我认识他们，我们是老相识了。”——“但是我讨厌他们，”K直言不讳，他把目光从两个助手移到村长身上，又从村长移回到助手身上，发现三个人的微笑一模一样，毫无区别。“既然你们两个到这儿来了，”他带着试探的口吻说，“那么就留下来帮着村长太太找一份文件吧，文件上写着‘土地测量员’，字底下画着蓝道道。”村长并无反对意见。就是说，不许K做的，两个助手倒可以做，两人也立即投身纸堆，甩开膀子大干起来，可是，与其说他们在找文件，倒不如说在纸堆里乱翻一气，而且，当一个人正读文件标题时，另一个人总是一把将文件从他手中抢走。那女人则跪在那现在已经空空如也的柜子前面，看样子根本没在继续找，至少现在蜡烛离她老远老远。



...



“您认识施瓦尔策吗？”K问。

“不认识，”村长说，“你大概认识吧，米齐？也不认识。不，我们不认识这个人。”

“这真是怪事，”K说，“他是城堡一位副主事的儿子。”

“亲爱的土地测量员先生，”村长说，“您怎么可以要求我认识所有城堡副主事的所有儿子？”

“很好，”K说，“那么您就只好相信我说的，承认他就是城堡一位副主事的儿子。同这个施瓦尔策，我在到这里的当天就有过一次不愉快的争论。后来他打电话给一个叫弗里茨的副主事询问，答复是我已经被聘任为土地测量员了。这您又怎么解释呢，村长先生？”

“很简单，”村长说，“这一点正好说明您从来还没有同我们的官府接触过嘛！所有这一切接触统统是表面文章，由于对情况完全无知，您竟将它们信以为真了。至于说到电话嘛，您瞧，我和官府在公务上的联系怎么说也是够多的了吧，可是我这里就没有电话。在酒店一类地方电话可能有相当大的好处，好比是八音盒那一类东西，然而更大的作用也就没有了。您在这里打过电话吗？打过？好，那您也许能明白我的意思。城堡里的电话显然非常灵；我听人说过，那里总是一个电话接着一个电话，这当然使工作的进度加快许多。那边连续不断的电话，我们这边电话机里听起来就是不停的嗡嗡声和歌唱声，这您一定也听到了。但我们这里的电话所能传达的唯一正确可信的东西，也就只是这种沙沙声和歌唱声罢了，别的信息全不是那么回事。我们同城堡之间并无确定的电话线路，没有总机接转我们打去的电话；从这里给城堡中某人挂电话时，那边最下一级的办事部门的电话就都一齐响起来，甚至于，这点我知道得很清楚，要不是城堡里几乎所有电话机上的响铃装置都关上了，那么所有的电话机就都会响起来的。偶尔有个别过于疲劳的官员觉着需要放松一下消遣消遣，特别是晚上或夜间，他就会把响铃装置打开；那时我们就能听到回话了，可那叫什么回话，不过是开开玩笑而已。这也是完全合情合理的。谁有权利为了想解决自己一点小小的私人困难而用电话铃声去干扰那些极为重要的、忙得人晕头转向的紧张工作？我也不明白为什么连一个外乡人也会以为，要是他比如说打电话给索尔蒂尼，那么接电话答复他的就一定是索尔蒂尼了。说多半会是另一个毫不相干的部门的一个小小记录员，这可能性不是反倒更大些吗？反过来说，在极少有的时候当然也会出现打电话给小记录员时反倒是索尔蒂尼本人来接电话的情况。那时当然是趁着还没听到他开腔说第一句话就赶紧跑开为妙。”
\par\vspace{0.3em}{\footnotesize\color{gray}（第五章）}
\newpage
\section{土地测量员聘用始末}
\vspace{-0.3em}{\color{gray}\small\kaifont \noindent K 不只要一个简单的说法。他要搞清楚，聘他来做土地测量员的这个决定到底是怎么跑出来的。村长讲了很久。故事的核心不是答案。}\vspace{0.5em}
\vspace{0.3em}

\noindent\rule{\textwidth}{0.4pt}

\vspace{0.3em}



“这么说这两个助手，”村长又开口了，洋洋得意地微笑着，似乎在说：所有这些全是他一手安排的，而谁都无法预想到这一点，“这么说，他们让您感到讨厌，可他们是您自己的助手呵。”——“不对，”K冷冷地说，“他们是我到此地后才跑到我身边来的。”——“怎么，跑来的，”村长说，“您的意思大概是说派来的吧。”——“好吧，就算是派来的吧，”K说。“但说他们是从天外飞来的也未尝不可，这种派法也够欠考虑的了。”——“这里办什么事都不会欠考虑，”村长说，急得连腿疼也忘了，一骨碌坐起来。“办什么事都不欠考虑，”K说，“那么招聘我这事又怎么说？”——“聘用您也是经过深思熟虑的，”村长说，“只是有一些次要因素起了干扰作用罢了，我可以用有关的文件向您证明这一点。”——“这些文件恐怕找不到了，”K说。“找不到？”村长叫起来。“米齐，请快点找！不过，没有文件我也可以先把事情的经过讲给您听听。刚才我提到的那份指令，我们接到后当即回复上头表示感谢，说我们不需要土地测量员。但是这封回信看来是没有送回发出指令的那个部去，我管它叫A部吧，而是误送到另一个部，送到B部去了。这样A部就一直没有得到回音，然而可惜连B部也没有收到我们复信的全部；一种可能是公文袋里装的信落在村公所，再一种可能是在半路上遗失了——肯定不会在部里遗失，这我敢担保，总之，B部也只收到一个空公文袋，袋上只注明封套内所装的、可惜实际上并不在内的文件是关于聘任土地测量员一事的。A部呢，一直还在等着我们的回复，这个部虽然已将此事记录在案，但部主管人指望着我们会答复，等着在收到答复之后他们或者聘任土地测量员，或者根据需要继续同我们就此事进行磋商。这种坐等下面回信的事绝非绝无仅有，这是可以理解的，在办什么事都讲求准确、一丝不苟的地方这样做也无可厚非。可是结果呢，因为他指望着我们，所以也就不大重视部里那份有关记录，弄到后来就把这事忘记了。但是在B部就不同，公文袋送到了一位以认真负责闻名的部主管人手里，他叫索尔蒂尼，是个意大利人；连我这个内行人也不明白为什么像他那样能干的人会被长期安置在那个可说很低的职位上。这位索尔蒂尼自然把空文件袋派人送回给我们要求补上内容。但是，从A部发出第一封函件到那时已经过了好几个月，甚至已经好几年了；这也不难理解，因为按常规如果一份文件送对路了，那么最迟在第二天也就送到部里，当天就可以处理；但如果没送对路——在组织非常严密完善的情况下，文件要走这条错误的路简直就必须费很大力气去找，否则是找不到的，那么，那么当然就会拖很长时间了。因此当我们收到索尔蒂尼的便函时，就只能隐隐约约地记起这件事来，当时我们管这事的只有两人，就是米齐和我，那位老师还没有分配给我，我们两个只能对最重要的事保存一份副本，简短地说吧，我们只能给他一个很笼统的回答，说我们完全不知道聘任一事，我们这里也不需要什么土地测量员。”



“但是，”村长顿住了，似乎发觉在乐此不疲的叙述中忘其所以，把话扯远了，或者至少有可能扯远了，“这个故事您听着不觉得腻味吧？”



“不，”K说，“我觉得挺有意思，挺解闷的。”



然后村长说：“我讲这些给您听不是为了给您解闷。”



“我所以觉得它有意思，”K说，“仅仅因为我听了这事后，对那种可笑的混乱算是窥见了一斑，这种混乱有时能决定一个人的命运哟。”



“您还没有真正窥见一斑，”村长正色说，“我可以接着跟您讲。我们的回答当然满足不了索尔蒂尼那样的人。我很佩服这个人，尽管一听到他的名字就头疼。为什么？因为他谁都不信赖，比如，即使他有过无数次机会进行了解，完全知道那是个最可靠的人，可到下一次再碰上时他还是不信任那个人，好像他根本不认识人家似的，或者更正确些说，好像他知道那人是个无赖。我觉得这是对的，一个公职人员就应该这样；可惜我的天性使我不能按这条原则办事，你不是看见了吗，我现在对您这个外乡人把什么底都亮出来了，我这叫做禀性难移，没法子。而索尔蒂尼一接到我们的回信就起了疑。于是便有了一连串的通信往来。索尔蒂尼问我为什么突然想起说不要聘用土地测量员；我仗着米齐的好记性回信说，这事最初是上头提出来的呀（至于事实上是另一个部发来的文件，这一点我们早忘记了）；索尔蒂尼对此的说法是：为什么我到现在才提起上级的这封公函；我又回复他说：因为我现在才想起这封公函来嘛；索尔蒂尼呢：这真是太奇怪了；我：对于拖了那么长时间的一件事，这是一点也不奇怪的；索尔蒂尼：这事确实很奇怪，因为我想起来的那封公函并不存在；我：当然不存在啦，因为关于这事的全部文件都丢失了；索尔蒂尼：如果确有那第一封函件，就必定会有一条有关的记录在，然而这样一条记录并不存在。这时我无话可说了，因为要说是索尔蒂尼那个部出了差错，这话我是既不敢讲也不相信的。您，土地测量员先生，也许会暗暗埋怨索尔蒂尼，心想只要他考虑一下我的话，至少会去向另外的一些部打听打听这事吧。然而真要这样做恰恰是不对的，我不想让您哪怕只是在心里产生一种索尔蒂尼身上还有点毛病的印象。压根不考虑有出错的可能性，正是衙门的一条工作原则。由于组织十分完备严密，这条原则是正确的，而如果想达到极高的办事效率，它也是必要的。所以说索尔蒂尼根本就不能去向其他各部打听，而且就算他去打听，这些部门也决不会答复他，因为它们马上就会感觉出：这是在追究一个出差错的可能性。”



“村长先生，请允许我打断您一下，向您提个问题，”K说，〔8〕“您不是提到过一个检查机关吗？按照您对这里办事情况的描述，如果设想一下这里对这些事居然会没有监督检查，这简直令人不堪忍受。”



“您非常严格，”村长说。“但您就是再严格一千倍，同衙门对待自己的那种严格相比，您的严格也仍然等于零。只有一个地地道道的外乡人才会提出您提的问题。要问有没有检查机关吗？这里有的只是检查机关。当然啰，这些检查机关不是为了查出一般意义上的错误而设立的，因为这里是不会出错的，即便某次出了一个错，如您的情况，但是谁敢下定论，说这肯定是一个错误！”



“这个说法真是太新鲜了！”K叫起来。



“对我来说这是很旧很旧的说法，”村长说。“我同您差不多，也确信是出了一个错，索尔蒂尼因为对这事伤透脑筋毫无进展而得了重病，第一批检查机关官员来了，亏得他们揭出了错误的来源，看出这里有一个错，可是谁敢拍着胸脯说，第二批检查机关官员也会做出同样的判断，第三批和以后各批也都会得出同样的结论？”



“就算是这样吧，”K说，“不过我最好还是不要对这些考虑妄加评论，我也是头一次听说这些检查机关，当然还不可能完全了解它们。但我觉得这里必须区别两点：第一是机关内部发生的事以及仍然从机关的角度出发对这些事可能有这样或那样的看法，第二是我作为一个具体的个人，我这个身在公职机关以外的人，现在正面临着公职机关无视我的存在的危险，这种无视现实是如此荒谬，以致我直到现在总不大相信这危险是实际存在着的。对于第一点，很可能您村长先生刚才讲的那些话都是对的，您对内情了解得如此了如指掌令人叹服，只是我在听了您这番话后也想听您说说我个人的事情了。”



“我也正想来谈这个，”村长说，“但如果我不预先说明几句，那您是听不懂我的话的，我刚才提起检查机关看来是为时过早了。所以现在我要再回过头来谈谈我跟索尔蒂尼的分歧。我已经说过，渐渐的我招架不住了。但索尔蒂尼哪怕只比他的对手略胜一筹，他就等于大获全胜了，因为这样一来他办起事来精力就更集中、精神就更饱满，就更加指挥若定；结果是他的手下败将一见他就胆战心惊，而他手下败将的敌人看着他就会拍案叫绝。只因为我对这后一点在另外一些事情上也有亲身体会，我才能像现在这样讲他这个人。不过我还从未有幸亲眼见过他，他不能到下面来，他真是日理万机，公务繁忙至极，有人对我描述过他的房间，说四壁全被大捆大捆的卷宗挡满，一捆摞一捆，形成了一根根高大的方柱，而且这些还仅仅是索尔蒂尼当时正在处理的文件，由于不断从一捆捆卷宗中抽出、插进文件，并且又都是为赶时间匆匆忙忙地做，这些卷宗堆成的柱子就不断倒塌下来，正是这种接二连三、每隔一会儿就出现一次的轰然巨响，成了索尔蒂尼办公室的突出特征。这也难怪，索尔蒂尼是个一丝不苟的工作人员，不论是最小的事还是最大的事，他都一视同仁，一样认真对待。”



“村长先生，”K说，“您总把我的事叫做小事一桩，但我这件事使许多公职人员忙个不亦乐乎，或许它在开始时只是一桩很小的事吧，可是由于像索尔蒂尼那样的官员们积极参与，它确实已经变成一件大事了。这一点我感到非常遗憾，这是完全违反我的本意的，我不敢有那样的野心：想让关于我的事情的卷宗也变成高大的柱子又随时哗啦啦塌下来，我只希望做个小小的土地测量员，在一张小小的绘图桌旁安安心心地工作就心满意足了。”



“不，”村长说，“您的事不是大事，在这点上您没有理由抱怨，您的事确实是诸多小事中最小的一件。工作量大小并不决定工作的重要性，如果您这样想那您对衙门的了解就太皮毛了。但即使工作量大小能说明问题是否重要，您这件事也仍然是一粒微不足道的芝麻，普遍的事由，即那些没有出所谓差错的事，办理起来工作量就大得多，工作的成果自然也丰富得多，不过您还根本不知道您的事给人造成的实际上的工作量呢，我现在就来说说这个。首先要说的是，索尔蒂尼虽然不再追问我，但他手下的工作人员来了，每天都要在贵宾楼对一些有声望的村民进行几次质询并作翔实记录。大部分人支持我的说法，然而也有少数人起了疑心；土地测量问题触及农民的根本利益，农民们嗅出有人在算计他们，觉得这里头有些不公道的事，他们还找到了一个领头的，于是索尔蒂尼从他们提供的情况中得出结论，确信：如果我当初把问题提到村民大会上讨论，是不会出现全体一致反对聘用土地测量员的情况的。〔9〕这样一来一件原本是明摆着谁都明白的小事——即不需要土地测量员——不管怎么说至少被弄得成为一个问题了。在这件事上，有一个叫布伦斯维克的表现得特别突出——您大概不认识他吧，他也许不是坏人，但又蠢又狂，他是拉塞曼的一个妹夫。”



“是那个鞣皮匠拉塞曼吗？”K问，接着便描述了他在拉塞曼家见到的那个大胡子的长相。



“是的，正是他，”村长说。



“我还认识他的女人，”K有点有一搭没一搭地说。



“这是有可能的，”村长说，然后便沉默了。



“她很美，就是气色不大好，病恹恹的。她大概是城堡里来的人吧？”这后一句话是带着疑问语气说出来的。



村长看了看表，将药水倒进一把汤匙里，急急忙忙灌下肚去。



“您大概只认得城堡里那些办事机构吧？”K颇不客气地问道。



“对，”村长带着自嘲然而却是感激的微笑说道。“机关也是最重要的呵。至于说到布伦斯维克：要是我们能把他从村里驱逐出去，几乎对每个人来说都是件大快人心的事，拉塞曼也不会是最不开心的一个。可是当时布伦斯维克给自己拉了一些支持者：他虽然不是演说家，倒也会大喊大叫，而这样对某些人也能起作用。于是，终归我还是被迫把这件事提交村民大会讨论。不过这也就是布伦斯维克取得的唯一的胜利了，因为到了村民大会上，当然是绝大多数人反对聘用土地测量员的。这次大会也又过去好几年了，但是这些年来这件事一直没有平息，这部分是因为索尔蒂尼太认真，他通过极为周密细致的调查，力图摸清多数派和反对派的动机，再就是因为布伦斯维克的愚蠢和野心，他同官府有各种各样的私人关系，而且有层出不穷的、异想天开的新招数去利用这些关系。当然索尔蒂尼不会上布伦斯维克的当，布伦斯维克怎么瞒得了索尔蒂尼？——可是，恰恰是为了不致上当受骗，需要进行新的调查，新的调查还没有结束，布伦斯维克又想出了新花招，他这人很灵活，这正是他那愚蠢的一种表现。好，现在我来谈谈我们政府机构的一个特殊性质。与它的精细严密相适应，它也极为敏感。如果一件事已经深思熟虑了很久很久，那么即使考虑还没有结束也会出现这种情况：在某个事先无法预见事后也无法找到的地方会闪电般突然冒出一个解决办法，这办法虽说多半是正确地了结了那件事，但无论如何总带有一定的随意性吧。情形就好像是：政府机关再也受不了那绷得太紧的弦、受不了同一件也许本身是微不足道的事情那长年累月的刺激，从而自发地、不等哪个官员来过问就作出决定了。这当然不是说出现了什么奇迹，可以肯定地说是某位官员挥笔写下了这个解决方案，或是作了一个口头的决定，但至少从我们这里，从此地，甚至从公职机关也怎么都不可能弄清楚究竟是哪位官员在这事上作了定夺，是根据什么理由这样决定的了。这些全是检查机关在很久以后才查明的；而我们呢，就永远不得而知了，实际上这事过了那么久恐怕也不会再有谁感兴趣了。我说过，正是这类决定，它们多半是正确的，是很好的，只有一点不足，就是——这类事通常都是如此——一般人知道得太晚，因此在知道前的一段长时间里一直还在对那早已决定了的事进行热烈的讨论。我不知道在您的事情上是否也发出过一个这样的决定。——有迹象说明发出过，也有迹象说明没有发出过；假如已经发出，那么聘任书就可能早已寄给了您，然后您就会长途跋涉到此地来，这些花费了很多时间，而与此同时索尔蒂尼却还会一直在这里为这同一件事操劳着，累得筋疲力尽，布伦斯维克也还会绞尽脑汁搞阴谋诡计，而我就会在他们两人中间受尽了夹板气。这个可能性只是我的揣测，我确切知道的是这一点：一个检查机关查出来好多年前A部曾经就聘任土地测量员一事向村里发出过询问函件，而后直至现在一直没有收到回音。新近又有官员向我查询这事，这样一来全部事实当然也就真相大白了，A部对我提供的不需要土地测量员这一回答表示满意，索尔蒂尼则不得不看到：原来他并不主管这件事情，并且还做了许多无用的、让人伤透脑筋、把人弄得精疲力竭的虚功。——当然，他是无辜的。要不是这段时间也同以往一样新的工作任务不断从四面八方源源而来，要不是您的事情说到底仅仅是桩很小的事——几乎可以说是诸多小事当中最小的一件，那么我们大家一定都会感到如释重负了，我相信索尔蒂尼自己也会这样，只有布伦斯维克有一肚子怨气，但那只会让人觉得可笑罢了。好，现在，恰恰是现在，在整个这件事总算善始善终顺利了结之后——从那时起到今天又已经过去好久，您突然又冒了出来，事情似乎又像要从头开始的样子，请您，土地测量员先生，设想一下我有多么失望呵。现在我下了大决心在我力所能及的范围内决不让这类事再发生，这一点想来您是很清楚也能理解的吧？”
\par\vspace{0.3em}{\footnotesize\color{gray}（第五章）}
\newpage
\section{巴纳巴斯送信}
\vspace{-0.3em}{\color{gray}\small\kaifont \noindent 巴纳巴斯是 K 的信使，也是我们在 Day 2 就认识的人。但直到今天，他的姐姐奥尔嘉才把他每天在城堡办公厅里经历的事情讲清楚。他没有工作服，所以他连自己算不算城堡的人都不知道。他穿梭在栅栏之间，那些允许他穿过的，和不允许他穿过的，外表一模一样。克拉姆给他派任务——听起来很荣耀——但信是秘书写的，内容是过时的。办公厅里的设置不是为了让人找到正确的路，是为了让外面的人找不到任何路。}\vspace{0.5em}
\vspace{0.3em}

\noindent\rule{\textwidth}{0.4pt}

\vspace{0.3em}



总而言之，巴纳巴斯绝对谈不上是职位较高的办事人员，那么照理说他就应该是个低级的办事人员了吧？可是那些低级办事人员全都有工作服，至少，在他们下到村里来时是穿着工作服的，这不是真正的制服，各人穿的也有很多小地方不一样，可是不管怎么说看到这种服装一眼便能认出是城堡来的人员，在贵宾楼你不是也见到过这样一些人嘛。这种服装最引人注目的一点就是多半都非常紧身，干农活的或者做手工的人恐怕是没法穿的。好了，巴纳巴斯没有这种衣服；这不光是让人感觉寒碜，觉得脸上无光，这倒还能忍受，严重的是让你对什么事都怀疑起来，特别是在情绪低落消沉时更是这样——我们，就是巴纳巴斯和我，我们有时会情绪很低落的，我说有时，并不是说这样的时候就很少很少。遇上这种时候我们就对什么都怀疑起来，我们琢磨：巴纳巴斯做的工作到底算不算是城堡的工作；不错，他是在那些办公厅里进进出出，但那些办公厅就是真正的城堡吗？就说城堡肯定有不少办公厅吧，但是，允许巴纳巴斯进出的那些办公厅是城堡的办公厅吗？他是走进了一些办公厅；可那毕竟只是所有办公厅中的一部分呀，那儿另有些栅栏门，栅栏门后面还有另外一些办公厅。人家倒也没有禁止他穿过这些栅栏，可他不可能再往前走了，因为在栅栏这边他就找到了他的那几个上司，他们听完他的报告就让他走人了。另外，在那里行动一直是受到监视的，至少到过那里的人有这个感觉。再说即使他穿过了栅栏又有什么用处呢，他在那边又没有要办的事，那样做岂不是成了一个擅自闯入公务重地的异己分子？你也不能把这些栅栏设想成某种不许越过的界线，巴纳巴斯也一再让我注意这一点。原来，在他进出的那些办公厅里也是有栅栏的；这就是说也有那么一些栅栏是他经常穿行的，而且这些栅栏跟他还没有穿越的那些看上去并没有两样，所以说也不能一开始就猜想在他没有穿越的栅栏后面会有跟他去过的那些办公厅大不相同的什么别的办公厅。不就是我们只在情绪低落消沉时才有这种想法嘛！就这样，我们的疑心越来越大，根本无法抗拒。再举个例子说吧，巴纳巴斯是跟一些官员说话，巴纳巴斯是得到人家让他传送的信，可那究竟是些什么官员，是些什么信呵！他自己说他现在是分配到克拉姆手下了，克拉姆本人亲自给他派任务。好，这应该说是很抬举他了吧，连一些高级管事也没有争取到这份差事，唔，恐怕这是把他抬得太高了，真让人担心会摔个粉身碎骨的。你想想看，直接分给克拉姆调遣，同他嘴对嘴地说话！真是这么回事吗？



...



你知道，弗丽达不怎么喜欢我，她是怎么也不乐意让我见到克拉姆的，但他的相貌在村里当然是人人都知道的，有少数人见过他，所有的人都听说过他，于是从见过的人得到的印象、道听途说听来的材料、再加上某些别有用心的人的添油加醋的描述，就勾画出了一幅克拉姆画像，这画像也许大体上符合实际，但只能说是大概像。除基本轮廓外，别的方面就有比较大的出入了，而这些描述的出入之处也许还没有克拉姆的真正外貌的变化来得大呢。据说他到村子里来时完全变了模样，而到离开村子时又变一个样，喝啤酒前是一个样，喝完后又是另一个样，醒时一个样，睡着了又一个样，独自一人时一个样，跟人谈话时又一个样，根据这些，说他在上头城堡里时完完全全是另一个人就由不得你不信了。再就是，即使那些有关他在村里时的外貌的说法，也有相当大的出入，关于他的个子、姿势、胖瘦、胡子等方面的说法都很不一致，只有在穿着一点上，幸好大家的说法总算一致了：都说他总是穿同一件衣服，一件有长摆的黑色上装。为什么会有这么些不一样的说法？当然不是克拉姆会变魔术，而是因为见到他的人当时的心情、激动的程度、各自抱的希望或绝望的复杂心境千差万别，数也数不清，加上见到他的人多半都只被允许匆匆一面，考虑到这些因素，有那么多不同说法就非常容易理解了。
\par\vspace{0.3em}{\footnotesize\color{gray}（第十五章）}
\end{document}